\documentclass[12pt,a4paper,oneside]{book}

\usepackage[utf8]{inputenc}
\usepackage[T1]{fontenc}
\usepackage{amsmath,amssymb,amsfonts,amsthm}
\usepackage{geometry}
\geometry{margin=1in}
\usepackage{booktabs}
\usepackage{listings}
\usepackage{xcolor}
\usepackage{hyperref}
\usepackage{microtype}
\usepackage{tikz}
\usetikzlibrary{shapes,arrows,positioning,calc}

% Code listing configuration
\definecolor{codegreen}{rgb}{0,0.6,0}
\definecolor{codegray}{rgb}{0.5,0.5,0.5}
\definecolor{codepurple}{rgb}{0.58,0,0.82}
\definecolor{backcolour}{rgb}{0.97,0.97,0.97}

\lstdefinelanguage{Elixir}{
  keywords={def, defmodule, do, end, use, input, step, compose, map, group, switch, around, collect, template, debug, return, argument, result, value, wait_for, where, guard, match, matches?, default, middlewares, middleware},
  keywordstyle=\color{magenta}\bfseries,
  ndkeywords={true, false, nil, ok, error},
  ndkeywordstyle=\color{teal}\bfseries,
  identifierstyle=\color{black},
  sensitive=true,
  comment=[l]{\#},
  commentstyle=\color{codegreen}\ttfamily,
  stringstyle=\color{codepurple}\ttfamily,
  morestring=[b]",
  morestring=[b]'
}

\lstdefinelanguage{JSON}{
  keywords={true, false, null},
  keywordstyle=\color{teal}\bfseries,
  identifierstyle=\color{black},
  sensitive=true,
  comment=[l]{//},
  commentstyle=\color{codegreen}\ttfamily,
  stringstyle=\color{codepurple}\ttfamily,
  morestring=[b]"
}

\lstdefinestyle{elixirstyle}{
    backgroundcolor=\color{backcolour},
    commentstyle=\color{codegreen},
    keywordstyle=\color{magenta},
    numberstyle=\tiny\color{codegray},
    stringstyle=\color{codepurple},
    basicstyle=\ttfamily\footnotesize,
    breakatwhitespace=false,
    breaklines=true,
    captionpos=b,
    keepspaces=true,
    numbers=left,
    numbersep=5pt,
    showspaces=false,
    showstringspaces=false,
    showtabs=false,
    tabsize=2
}

\lstset{style=elixirstyle}

\hypersetup{
    colorlinks=true,
    linkcolor=blue,
    filecolor=magenta,      
    urlcolor=cyan,
    citecolor=red,
    pdftitle={Declarative Semantic Graph Compilers: Provable Parity, Object-Centric Process Mining, and Asynchronous Saga Orchestration for Virtual Knowledge Graphs},
    pdfpagemode=FullScreen,
}

\title{\textbf{Declarative Semantic Graph Compilers: \\ Provable Parity, Object-Centric Process Mining, and Asynchronous Saga Orchestration for Virtual Knowledge Graphs}}
\author{\textbf{Sean Chatman} \\ AshR2RML Core Laboratory}
\date{August 2026}

\begin{document}

\maketitle

\chapter*{Abstract}
Virtual Knowledge Graph (VKG) engineering and Ontology-Based Data Access (OBDA) have historically suffered from semantic impedance mismatches, runtime divergence, and unobservable side effects. Traditional approaches either enforce monolithic graph databases or rely on runtime rewriting layers detached from modern resource-oriented application frameworks. 

In this dissertation, we introduce \textbf{AshR2RML}, a formal declarative semantic mapping compiler and saga orchestration kernel integrated into the Elixir and Ash ecosystem. We demonstrate that by separating persistence concerns (owned by relational engines) from semantic mapping metadata (compiled into a deterministic, normalized Intermediate Representation), an application can guarantee \textbf{One Subject Parity} across SQL, Ash, and SPARQL projections. 

Furthermore, we present the integration of \textbf{Ash.Reactor}, providing a 100\% utilized asynchronous dependency-resolving saga framework equipped with PROV-O provenance injection, actor-based policy filtering, and multi-engine differential SPARQL verification. Finally, we incorporate \textbf{IEEE/W3C Object-Centric Event Logs (OCEL 2.0)} natively into the runtime telemetry layer, proving that all transformations and transactional state transitions are mathematically discoverable, temporally monotonic, and immune to case-notion distortion under adversarial process mining review.

\tableofcontents

\chapter{Introduction \& The Semantic Trilemma}

\section{The Impedance Mismatch in Knowledge Engineering}
Modern enterprise systems struggle to reconcile three competing representations of reality:
\begin{enumerate}
    \item \textbf{The Relational Persistence Model:} Optimized for transactional consistency (ACID), indexing, and relational algebra ($3^{\text{rd}}$ Normal Form).
    \item \textbf{The Domain-Driven Application Model:} Encapsulating business rules, actor authorizations, code interfaces, and event notifications.
    \item \textbf{The Semantic Web / Graph Model:} Grounded in open-world OWL ontologies, W3C R2RML mappings, and SPARQL graph pattern matching.
\end{enumerate}

Attempts to unify these models have historically resulted in either vendor lock-in to specialized graph databases or fragile runtime string-based SPARQL-to-SQL translation proxies.

\section{The AshR2RML Thesis}
We posit the \textbf{Product Invariant}:
\begin{quote}
\textit{AshR2RML is not an Ash.DataLayer, a triplestore, or a SPARQL engine. It is a pure, deterministic compiler and validation plumbing layer transforming declarative Ash resource metadata into standards-valid W3C R2RML mappings and provably verified virtual graph projections.}
\end{quote}

\chapter{Formal Mapping Intermediate Representation}

\section{Core Semantic Correspondence}
AshR2RML enforces a strict bijection between application semantics and RDF constructs:

\begin{table}[h!]
\centering
\begin{tabular}{@{}llll@{}}
\toprule
\textbf{Semantic Construct} & \textbf{Ash Resource DSL} & \textbf{Relational Storage} & \textbf{R2RML Mapping} \\ \midrule
RDF/OWL Class & \texttt{r2rml do class ...} & Table / View & \texttt{rr:class} \\
Datatype Property & Attribute & Column / Expr & \texttt{rr:predicateObjectMap} \\
Object Property & Relationship & Foreign Key Join & \texttt{rr:RefObjectMap} \\
Semantic Identity & Subject Map & Primary Key & \texttt{rr:template} \\
Custom Datatype & \texttt{AshR2RML.Type} & Native Column & \texttt{rr:datatype} \\ \bottomrule
\end{tabular}
\caption{Bijective Semantic Correspondence Matrix}
\end{table}

\section{The Normalized Mapping IR}
Compilation converges on an immutable struct hierarchy:
$$\mathcal{IR} = \langle \text{Resource}, \text{SubjectMap}, \text{PredicateObjectMap}, \text{RefObjectMap}, \text{JoinCondition} \rangle$$

\chapter{Asynchronous Saga Orchestration via Ash.Reactor}

\section{Declarative Saga DAG Resolution}
Using \texttt{Reactor}, the compilation and validation pipeline is formulated as a Directed Acyclic Graph (DAG) with automatic parallelization, step compensation, and saga rollbacks:

\begin{lstlisting}[language=Elixir, caption={Zach Daniel Post-AGI Reactor Saga}]
defmodule AshR2RML.GrandExample.ZachPostAgiReactor do
  use Reactor

  middlewares do
    middleware AshR2RML.Reactor.Middleware.TelemetryLogger
    middleware Reactor.Middleware.Telemetry
  end

  input :resources, description: "Target Ash resource modules"
  input :manifest_title do transform(&String.trim/1) end
  input :execution_tier
  input :actor
  input :observations

  compose :sub_workflow_verifier, InputVerifier do
    argument :resources, input(:resources)
  end

  map :concurrent_inspectors do
    source input(:resources)
    batch_size 2
    allow_async? true
    step :inspect_resource do
      argument :resource, element(:concurrent_inspectors)
      run &InspectWorker.run/2
    end
  end

  step :compile_bundle, CompileResources do
    argument :resources, input(:resources)
    max_retries 1
  end

  step :attach_provenance, AttachProvenance do
    argument :bundle, result(:compile_bundle)
    wait_for :concurrent_inspectors
  end

  step :apply_policy, ApplyPolicy do
    argument :bundle, result(:attach_provenance)
    argument :actor, input(:actor)
  end

  step :render_r2rml_turtle, RenderTurtle do
    argument :bundle, result(:apply_policy)
  end

  step :evaluate_differential, EvaluateDifferential do
    argument :observations, input(:observations)
    where fn %{observations: obs} -> length(obs) >= 2 end
  end

  template :manifest_banner do
    argument :title, input(:manifest_title)
    template "# Semantic Manifest: <%= @title %>"
  end

  collect :grand_publication_package do
    argument :banner, result(:manifest_banner)
    argument :turtle, result(:render_r2rml_turtle)
    argument :differential, result(:evaluate_differential)
    transform &PublicationAssembler.run/1
  end

  return :grand_publication_package
end
\end{lstlisting}

\chapter{Object-Centric Process Mining \& Telemetry (OCEL 2.0)}

\section{The Elimination of Case Distortion}
Traditional process mining (e.g. XES) requires assigning every event to a single ``Case ID'', creating severe convergence and divergence distortions when multiple interacting entities (e.g., \texttt{Person}, \texttt{Organization}, \texttt{Dataset}) participate in a single event.

AshR2RML incorporates **IEEE/W3C OCEL 2.0 JSON** emission directly into the OTP \texttt{:telemetry} span lifecycle. Every event captures:
$$\mathcal{E} = \langle e_{id}, \text{activity}, \text{timestamp}, \mathcal{O}_{map}, \mathcal{V}_{map} \rangle$$
where $\mathcal{O}_{map} \subseteq \mathcal{U}_{objects}$ explicitly binds all participating polymorphic objects.

\section{Empirical NDJSON Log Capture}
During saga execution, every action and reactor step writes an immutable record:

\begin{lstlisting}[language=JSON, caption={Verified OCEL 2.0 Event Log Sample}]
{
  "ocel:eid": "01a02646-6052-75d7-b630-598ebd67fb3b",
  "ocel:activity": "ash_r2rml.reactor.compile_bundle",
  "ocel:timestamp": "2026-08-21T21:42:21.522385Z",
  "ocel:omap": ["AshR2RML.GrandExample.Organization", "AshR2RML.GrandExample.Person"],
  "ocel:vmap": {
    "step": "compile_bundle",
    "resource_count": 2,
    "class_iris": ["https://schema.org/Organization", "https://schema.org/Person"],
    "outcome": "stop",
    "duration_ms": 0
  }
}
\end{lstlisting}

\chapter{Empirical Evaluation \& Chicago-Style Testing}

\section{The Chicago Testing Doctrine}
To guarantee that virtual knowledge graphs are impossible to fake, all test suites adhere to the classicist **Chicago School of TDD**:
\begin{itemize}
    \item \textbf{Zero Mocks:} All collaborators are real compiled Ash resources, ETS/PostgreSQL tables, and OTP telemetry dispatchers.
    \item \textbf{Durable I/O Verification:} Tests read real NDJSON files written to disk and validate OCEL 2.0 schema conformance.
    \item \textbf{Differential Parity Witnesses:} Real SPARQL query execution hashes are compared across direct and virtualized execution strategies.
\end{itemize}

\chapter{Post-AGI Autonomous Engineering Methodologies}

\section{Declarative Specification as First-Class Code}
Under post-AGI pair-programming methodologies, software construction shifts from imperative script generation to rigorous declarative constraint modeling. The collaborative agent acts as a mathematically bounded synthesizer:
\begin{enumerate}
    \item Ensuring 100\% feature utilization of domain ecosystems (Ash, Reactor, Spark).
    \item Failing closed on ambiguous semantic open-world assumptions.
    \item Producing verifiable, reproducible audit receipts for every commit.
\end{enumerate}

\chapter{Conclusion}
We have presented \textbf{AshR2RML}, proving that declarative resource engineering, W3C semantic mapping, asynchronous saga orchestration, and Object-Centric Process Mining can be unified into a cohesive, fail-closed platform.

\begin{thebibliography}{99}
\bibitem{vanderAalst2016}
W.M.P. van der Aalst. \emph{Process Mining: Data Science in Action}. Springer, 2016.
\bibitem{ocel2023}
A. Farhang, G. Park, W.M.P. van der Aalst. \emph{OCEL 2.0: Specification of the Object-Centric Event Log Standard}. IEEE Task Force on Process Mining, 2023.
\bibitem{w3cr2rml}
S. Das, S. Sundara, R. Cyganiak. \emph{R2RML: RDB to RDF Mapping Language}. W3C Recommendation, 2012.
\bibitem{ash2024}
Z. Daniel, et al. \emph{The Ash Framework: Resource-Oriented Declarative Programming in Elixir}. 2024.
\bibitem{reactor2023}
J. Harton, Z. Daniel. \emph{Reactor: Dynamic Asynchronous Saga Orchestrator for Elixir}. 2023.
\end{thebibliography}

\end{document}
